%=====================================================================
%  ch05 —— 原文 1208–1793：§3.3 Equivariance and Harish-Chandra modules
%=====================================================================

\subsection{等变性与 Harish-Chandra 模}

\ssszh 专家们立刻看出，局部化定理所承载的结构远不止范畴的等价。对具体的表示论应用
而言，主要的相容性如下。固定代数子群 $H \subset G$，为简单起见设它连通。

\ssszh 先从李理论一侧开始。与 $H$ 相伴的是 Harish-Chandra 模构成的全子范畴
%
\[
(\fg, H)\mod^\heartsuit \subset \fg\mod^\heartsuit.
\]
%
具体地说，一个 $\fg$-模是 Harish-Chandra 模，是指李代数 $\fh \subset \fg$ 的作用可积分
为 $H$ 的作用。由于我们假设 $H$ 连通，这是一条性质，而不是额外的结构。

限制到中心特征平凡的模，得到类似的全子范畴
%
\[
(\fg, H)\mod_0^\heartsuit \subset \fg\mod_0^\heartsuit.
\]
%

\ssszh 下面举几个简短的例子。

\begin{ex}
若 $\fg$ 是半单的，则中心商 $U_0(\fg)$ 只是平凡子群的 Harish-Chandra 模。
\end{ex}

\begin{ex} 平凡模 $k$ 是 $G$ 自身的 Harish-Chandra 模，从而也是任意子群的
Harish-Chandra 模。
\end{ex}

\begin{ex} Verma 模 $U(\fg) \underset{U(\fb)} \otimes \det(\fg/\fb)$ 是 $(\fg, B)$-模。
更一般地，对 $H$ 的任意表示 $W$，诱导模
%
\[
U(\fg) \underset{U(\fh)} \otimes W
\]
%
是 $(\fg, H)$-模，而每个 $(\fg, H)$-模都是这种模的商。
\end{ex}

\ssszh \label{s:eqdmod}下面转向几何一侧。与 $H$ 相伴的是 $H$-等变 D-模构成的全子范畴
%
\[
\Dmod(\sB)^{H, \heartsuit} \subset \Dmod(\sB)^\heartsuit.
\]
%

我们在任意带 $H$ 作用的代数簇 $X$ 的一般情形下回顾它的定义。

\begin{re} 对我们来说，代数簇按定义是 $k$ 上有限型的、拟紧、拟分离、既约的概形。
此外，本书对代数簇陈述的一切，对 $k$ 上有限型的拟紧拟分离概形同样成立，也就是说，
既约性并不起任何作用。读者若跳过这些技术细节，只考虑例如拟射影簇，也不会有多少损失。
\end{re}

初看之下，一个 D-模 $M$ 是 $H$-等变的，大致是说它沿 $X$ 中的 $H$-轨道为常值。
要得到真正的定义，注意我们因此应当要求：对 $H$ 的任意点 $h$ 与 $X$ 的任意点 $x$，
存在茎的同构
%
\begin{equation}\label{e:stalks}
i_x^! M \simeq i_{hx}^! M.
\end{equation}
%
当然，由于我们的空间并不离散，需要把这样的陈述在 $h$ 与 $x$ 变动时统一成一个族，
并与平行移动相容。为此，考虑作用映射与投影映射
%
\begin{equation}
\alpha, \pi: H \times X \rightarrow X,\quad \quad \alpha(h, x) = hx,
\quad \pi(h, x) = x.
\end{equation}
%

有了这些，$\Dmod(X)^\heartsuit$ 的对象 $M$ 是 $H$-等变的，当且仅当存在同构
%
\begin{equation} \label{e:wholething}
\alpha^! M \simeq \pi^! M.
\end{equation}
%
由于我们只要求这样的同构存在，这显然是 D-模的一条性质，而非额外的结构。

与函数的情形作比较是有启发的。即，固定一个作用在有限集上的有限群，不妨暗示性地
记它们为 $H(\Fq)$ 与 $X(\Fq)$。注意，$\Fun(X(\Fq))$ 中的函数 $m$ 是 $H(\Fq)$-不变的，
即沿轨道为常值，当且仅当它满足等式
%
\[
\alpha^! m = \pi^! m,
\]
%
其中 $\alpha, \pi: H(\Fq) \times X(\Fq) \rightarrow X(\Fq)$ 如上定义，而函数的
$!$-拉回如第 \ref{ss:opsfinitesets} 小节所述。

\ssszh 下面是一对基本而重要的等变 D-模的例子。

\begin{ex} 设 $X$ 由单个 $H$-轨道构成，即选定基点后形如 $H/K$，其中 $K$ 是 $H$ 的
某个子群。此时我们断言，$X$ 上的 $H$-等变 D-模与 $K$ 的连通分支群的表示之间，有一个
典则等价，即
%
\[
\Dmod(X)^{H, \heartsuit} \simeq \on{Rep}(\pi_0(K))^{\heartsuit}.
\]
%
事实上，容易看出 $X$ 上任一等变 D-模必是局部系统，而条件 \eqref{e:wholething}
等价于说沿 $H \rightarrow H/K$ 的拉回可平凡化。另一方面，纤维化
$K \rightarrow H \rightarrow H/K$ 所联系的同伦群长正合列在低次给出正合列
%
\[
\pi_1(H) \rightarrow \pi_1(H/K) \rightarrow \pi_0(K) \rightarrow 1,
\]
%
于是断言得证。 \label{e:oneorbit}
\end{ex}

\begin{ex} 设 $X$ 由有限多条 $H$-轨道构成。此时 $X$ 上的 $H$-等变 D-模范畴可由
各条轨道通过重粘（recollement）「拼接」而成。

其细节对我们并不那么重要。不过要提到下面这条基本推论：$X$ 上不可约的 $H$-等变
D-模，与各条轨道上不可约的 $H$-等变局部系统，处于典则的双射。具体地说，
$\Dmod(X)^{H, \heartsuit}$ 的每个不可约对象都支撑在单条轨道的闭包上，它在该轨道上的
限制是不可约的 $H$-等变局部系统。这样，由上一个例子，单对象的分类就归结为有限群论的
一道习题。

\label{e:bigbrain}
\end{ex}

\ssszh 两侧都已介绍，现在可以陈述所许诺的相容性了：局部化把 Harish-Chandra 模与
等变 D-模互换。

\begin{theo}[Beilinson--Bernstein, \cite{beilinson-bernstein81}]\label{t:heqv}
有交换图
%
\[
\xymatrix{\Dmod(\sB)^{H, \heartsuit} \ar[r]^\sim \ar[d] & (\fg,
H)\mod_0^\heartsuit \ar[d]  \\  \Dmod(\sB)^\heartsuit \ar[r]^\sim &
\fg\mod_0^\heartsuit,}
\]
%
其中竖直箭头是重言的包含函子，底部的水平箭头是局部化。
\end{theo}

\ssszh 事实上，定理 \ref{t:heqv} 在作了必要的调整（mutatis mutandis）之后，对任意的
$H$（不必连通）成立，对适当定义的等变 D-模与 Harish-Chandra 模的导出范畴也成立。
后一断言并不能由前一断言推出，因为这些导出范畴通常并不是其心（heart）的导出范畴，
而是稍为精细的对象。

\begin{re} 建议读者把上述关于导出范畴的微妙之处，看成如下形式的陈述：「要得到正确的答案，就必须稍微聪明一点。」然而，不难看出为什么朴素地用等变层的阿贝尔范畴的导出范畴
去处理会得到错误的答案。这里的定向原则是：$X$ 上的 $H$-等变 D-模，应当就是商叠
$X/H$ 上的 D-模，参看下面的注 \ref{r:introstacks}，因此应当与它的拓扑保持预期的关系。
在 $X = \on{pt}$ 这个最简单的情形，我们尤其应当希望
%
\[
\on{Hom}_{\Dmod(\on{pt})^H}(k, k) \simeq \on{C}^*(\on{pt}/H, k),
\]
%
即：常值层的导出自同态应当是分类空间的上同调。但既然在阿贝尔范畴的层次上
%
\[
\Dmod(\on{pt})^{H, \heartsuit} \simeq \on{Rep}(\pi_0(H)),
\]
%
我们单从阿贝尔范畴本身根本得不到任何高阶 Ext。这一修正归功于 Bernstein--Lunts，
我们将在第 \ref{s:catrepgrp} 节遇到它。
\end{re}

\begin{re} \label{r:introstacks}对不熟悉叠的读者，下面这个粗略的描述已足以应付我们的
需要。代数闭域 $k$ 上的 Artin 叠 $Y$，大致说来就像代数簇，但多了一条性质：$Y$ 的每个
$k$-点 $y$ 都带有一群「内部对称性。」

对我们最相关的一类 Artin 叠来自如下构造。给定代数群 $H$ 作用在代数簇 $X$ 上，可以
构造商叠 $X/H$。$X/H$ 的点由 $X$ 上的 $H$-轨道给出；而对 $X$ 的 $k$-点 $x$，
$X/H$ 中对应点的内部对称群恰是 $x$ 在 $H$ 中的稳定化子 $H_x$。也就是说，这个叠的
$k$-点并不是代数簇 $\pt$，即 $\Spec k$，而是它的商 $\pt/H_x$。一般地，Artin 叠的
$k$-点 $y$ 是某个代数群 $H_y$ 的分类空间 $\pt/H_y$，这个 $H_y$ 就是上面提到的内部
对称群。

作为简单的例子，考虑 $X = \mathbb{A}^1$ 带上由对合 $z \mapsto -z$ 诱导的
$H = \mathbb{Z}/2$ 的作用。若挖掉原点，这就限制为 $H$ 在 $\mathbb{G}_m$ 上的作用。
此时每条 $H$-轨道都是自由的，特别地，商中的点应当像普通代数簇那样具有平凡的内部
对称性。事实上，商 $\mathbb{G}_m/H$ 就是代数簇 $\mathbb{G}_m$，其中的商映射
%
\[
\mathbb{G}_m \rightarrow \mathbb{G}_m / H
\]
%
等同于平方映射 $\mathbb{G}_m \rightarrow \mathbb{G}_m$，$w \mapsto w^2$。
然而，完整的商 $\mathbb{A}^1/H$ 不是代数簇，因为对应原点的那个点已「轨道化」成了一份
$\pt/H$。直观地说，对一般的商叠，内部对称群记录了在从 $X$ 构造轨道空间 $X/H$ 时
发生的折叠。

最后要提到，与代数簇的情形相比，叠上的层的基本新特征是其纤维带有内部对称群的作用。
即，对代数群 $H$，有
%
\[
\QCoh(\pt/H) \simeq \on{Rep}(H) \quad \text{与} \quad \Dmod(\pt/H) \simeq
\on{C}_*(H)\mod,
\]
%
其中 $\on{C}_*(H)$ 记 $H$ 的 de Rham 上同调的对偶，并看作卷积下的 dg 代数。

因此，对 Artin 叠 $Y$ 上的凝聚层，它在 $k$-点 $y$ 处的 $*$-纤维带有 $H_y$ 的作用。
类似地，对 $Y$ 上的 D-模，它在 $y$ 处的 $!$-纤维带有 $\on{C}_*(H_y)$ 的作用，可以把
它想成 $H_y$ 的局部常值作用。特别地，该 $!$-纤维的上同调群带有连通分支群
$\pi_0(H_y)$ 的作用。
\end{re}

\ssszh 下面给出定理 \ref{t:heqv} 如何应用的一个基本例子。

由 \ref{t:heqv} 可知，$(\fg, H)\mod_0$ 中只有有限多个不可约对象，当且仅当 $\sB$ 上
只有有限多条 $H$-轨道。而且此时可以像例 \ref{e:bigbrain} 那样显式列举单对象。

作为具体例子，记 $T \subset GL_2$ 为极大环面，我们来分类
%
\[
(\fgl_2, T)\mod_0^\heartsuit
\]
%
的单对象。为此注意 $T$ 在 $\mathbb{P}^1$ 上有三条轨道——北极 $B_0$、南极
$B_\infty$，以及其余一切。由于 $T$ 在这三条轨道上的稳定化子都连通，可知共有三个
不可约对象。具体地说，它们是最高权为 $-2$ 的单 Verma 模、最低权为 $2$ 的单 Verma 模，
以及平凡模，即
%
\[
      U(\fgl_2) \underset{U(\fb_0)} \otimes \det(\fgl_2/\fb_0), \quad
      U(\fgl_2) \underset{U(\fb_\infty)} \otimes
      \det(\fgl_2/\fb_\infty), \quad \text{与} \quad k.
\]
%
还要注意，由类似的推理，当 $n > 2$ 时，$(\fgl_n, T)\mod_0^\heartsuit$ 中有无穷多个
单对象。

这段讨论相当软性，但我们要强调：更精细的问题，例如确定单模的特征，可以从几何入手，
而且通常翻译为关于轨道闭包拓扑的问题。

\ssszh 最后，就定理 \ref{t:heqv} 为何成立作两点定向性的说明。在下面的普遍性中讨论
会更清楚。\label{s:kandx}

设代数群 $K$ 作用在光滑代数簇 $X$ 上。如前面那样，若用 $\fk$ 记 $K$ 的李代数，
则有整体截面函子
%
\[
   \Gamma(X, -): \Dmod(X)^\heartsuit \rightarrow \fk\mod^\heartsuit.
\]
%
它有左伴随，由局部化给出
%
\[
   \on{Loc}: \fk\mod^\heartsuit \rightarrow \Dmod(X)^\heartsuit, \quad
   \quad M \mapsto \sD_X \underset{U(\fk) } \otimes M.
\]
%
我们断言，对任意代数子群 $H \subset K$，它限制为一个伴随对
%
\begin{equation} \label{e:locadj}
  \on{Loc}:(\fk, H)\mod^\heartsuit \leftrightarrows
  \Dmod(X)^{H, \heartsuit}: \Gamma(X, -).
\end{equation}
%
特别地，取 $K = G$、$X = \sB$，就重新得到定理 \ref{t:heqv}。

\ssszh 要看出 \eqref{e:locadj} 为何成立，需要先重新审视等变 D-模的定义。我们沿用
第 \ref{s:eqdmod} 小节的记号。对一般的群 $H$（不必连通），等变 D-模定义中的部分数据
是一个指定的同构
%
\[
\tau:  \alpha^! M \simeq \pi^! M,
\]
%
它满足单位性，并在 $H \times H \times X$ 上满足上圈条件（即结合性条件）。在茎的
层次上，上圈条件要求：把 $i^!_{h_1 h_2 x} M \simeq i_x^! M$ 认同为 $i_x^! M$ 的
两种方式——或对配对 $(h_1h_2, x)$ 用一次 $\tau$，或先对 $(h_2, x)$、再对
$(h_1, h_2x)$ 用两次 $\tau$——彼此一致。有一个很好的事实：若 $H$ 连通，则这样的
数据只要存在就唯一，而其存在性等价于先验地更弱的陈述 \eqref{e:wholething}。

\begin{re}
对不连通的群，等变性确实是为 D-模附加的进一步结构。例如可以验证：若 $X = \on{pt}$，
则 $\Dmod(\on{pt})^\heartsuit \simeq \Vect^\heartsuit$ 的对象 $W$ 上的
$H$-等变数据，等同于连通分支群 $\pi_0(H)$ 在 $W$ 上的表示。

读者或许从例 \ref{e:oneorbit} 中认出了这类表示。这不是巧合，事实上我们有
%
\[ \Dmod(H/K)^{H, \heartsuit} \simeq \Dmod(H)^{H \times K, \heartsuit}
\simeq \Dmod(\on{pt})^{K, \heartsuit}.  \]
%
用叠上的 D-模来说，这来自恒等式 $H \backslash H/K \simeq \on{pt}/K$。
\end{re}

\label{s:rigideq}

\ssszh \label{s:rigimon}第 \ref{s:rigideq} 小节中这种刚性化的等变性定义，或许会让
读者想起代数 $H$-模（即 $\sO_H$-余模）的定义。后者是一个向量空间 $W$，配有满足
单位性与上圈条件的余作用映射
%
\[
   W \rightarrow \sO_H \otimes W.
\]
%

\begin{re}若读者此前未见过这个定义，在目前的情形下它归结为下面这个或许更熟悉的描述：
$W$ 带有抽象群 $H(k)$ 的作用，是有限维表示 $M_\alpha$ 的并，并且关于这种子表示
$M_\alpha$ 的任意一组基，$H(k)$ 作用的矩阵系数都是 $H$ 上的正则函数。
\end{re}

把 $W$ 看作 $\on{pt}$ 上的拟凝聚层，这可以改述为 $\on{pt}$ 上一个 $H$-等变拟凝聚层
的数据，即沿 $\on{pt} \rightarrow \on{pt}/H$ 的下降数据。

有了这些，我们可以把 $H$-簇 $X$ 上 D-模的 $H$-等变数据重新表述如下。先设 $X$ 仿射，
并对 D-模 $M$ 用同一个字母记相应的 $D_X$-模，即其整体截面。

于是，$M$ 上的 $H$-等变数据，等同于在其底向量空间上给定 $H$ 的代数作用，满足
(i) 它与 $H$ 在 $D_X$ 上的作用相容（即从 $D_X$-模升级为 $H \ltimes D_X$-模），
(ii) 由此诱导的 $\fh$ 在 $M$ 上的无穷小作用，与用向量场的作用一致，即与下述复合一致
%
\begin{equation} \label{e:hch}
\fh \rightarrow T_X \rightarrow D_X.
\end{equation}
%
也就是说，人们把 $M$ 上由联络给出的 $\fh$ 作用，积分成群 $H$ 的作用。

对一般的 $X$（不必仿射），可以给出类似的定义。具体地说，D-模的等变数据是
(i) 其底拟凝聚层的等变数据，它与代数层 $\sD_X$ 上典则的等变数据相容，并满足上面
条件 (ii) 的类似物。

\ssszh 由这个表述可以看出，整体截面限制为一个函子
%
\[
         \Dmod(X)^{H, \heartsuit} \rightarrow (\fk, H)\mod^\heartsuit.
\]
%

我们再来勾勒局部化为何限制为一个函子
%
\[
   (\fk, H)\mod^\heartsuit \rightarrow \Dmod(X)^{H, \heartsuit}.
\]
%
出发点是：$\sD_X$ 自然地是 $H$-等变的拟凝聚代数层——这只是在说，函数与向量场可以
沿 $H$ 作用移动。然而，$\sD_X$ 并\emph{不}是 $H$-等变的 D-模。问题出在条件 (ii)，
即 $\fh$ 的两个无穷小作用必须一致。

要看出这一点，注意由 $\fh$ 在 $\sD_X$ 上的左作用所诱导的 $\fh$ 作用，形如
$h \otimes P \mapsto h \cdot P$，其中 $h$ 是 $\fh$ 的元素，$P$ 是 $\sD_X$ 的局部
截面。另一方面，由 $H$-等变性得到的 $\sD_X$ 上的无穷小 $\fh$ 作用必须由该代数的
导子给出，因此不能与左作用一致。于是它只能是伴随作用，即形如
%
\[
     h \otimes P \mapsto h\cdot P - P\cdot h.
\]
%

既然问题出在右乘的出现，读者就不会惊讶：用右乘与一个 $H$-可积模作张量应当有帮助。
确实，给定 $(\fk, H)$-模 $M$，容易看出
%
\begin{equation*}
\on{Loc}(M) = \sD_X \underset{U(\fk)} \otimes M
\end{equation*}
%
的底拟凝聚层，从 $H$ 在其张量因子上的作用继承了一份 $H$-等变数据。至于条件 (ii)，
注意这个作用的微分形如
%
\[
          h \otimes (P \otimes m) \mapsto h \cdot P \otimes m - P \cdot h
          \otimes m + P \otimes h \cdot m,
\]
%
其中 $h$ 与 $P$ 如前，$m$ 是 $M$ 的元素。但现在第二、三项相消，因为我们是在
$U(\fk)$ 上作张量，于是只剩下所希望的第一项。
